% =============================================================================
% 2. INTRODUCCIÓN
% =============================================================================

\section{Introducción}

\subsection{Escenario}

\subsubsection{Contexto del negocio}

Smart-Park es una plataforma web de tipo marketplace orientada a resolver la problemática de búsqueda y gestión de espacios de estacionamiento en la ciudad de Ayacucho (Huamanga), con capacidad de expansión a otras ciudades. Su modelo de negocio es un marketplace de varios inquilinos (\textit{multi-tenant}) bajo la modalidad SaaS para los estacionamientos afiliados.

El modelo consiste en conectar a los conductores con una red de estacionamientos afiliados. Cada local conserva autonomía operativa sobre sus espacios, su plano digital, sus tarifas por tipo de vehículo (por hora o por minuto, con recargo nocturno opcional), su tiempo de tolerancia, sus políticas de reserva, sus cámaras, su personal y sus operaciones de garita. La plataforma central se encarga de lo común a toda la red: la afiliación de nuevos locales, la seguridad y el control de accesos, la gestión financiera y las comisiones, la configuración global y la auditoría de los eventos del sistema.

Las fuentes de ingreso de la plataforma son cuatro: una comisión por transacción (del 10\,\% al 12\,\%), la suscripción mensual de los locales, la tarifa de servicio por reserva y los abonos mensuales.

\subsubsection{Problemática identificada en conductores}

Se aplicó una encuesta a 108 conductores de la ciudad, cuyos resultados sustentan la necesidad del sistema:

\begin{table}[H]
\caption{Resultados de Encuesta a Conductores}
\begin{tabularx}{\textwidth}{X r}
\toprule
\textbf{Indicador} & \textbf{Porcentaje} \\
\midrule
No siempre sabe, antes de dirigirse a un estacionamiento, si existen espacios disponibles & 90,7\,\% \\
Considera la congestión vehicular un problema ligado a la búsqueda de estacionamiento & 75,0\,\% \\
Señala que demora demasiado tiempo en encontrar un espacio libre & 75,9\,\% \\
Ha dejado de realizar una compra, trámite o actividad por no encontrar estacionamiento & 76,9\,\% \\
No conoce las tarifas antes de ingresar & 64,8\,\% \\
No puede reservar un espacio con anticipación & 66,7\,\% \\
Utilizaría una aplicación que muestre estacionamientos disponibles en tiempo real & 96,3\,\% \\
Utilizaría una aplicación para encontrar y reservar estacionamientos & 99,1\,\% \\
Utilizaría la solución si estuviera disponible & 98,1\,\% \\
Estaría dispuesto a pagar un pequeño monto de garantía para reservar un espacio & 98,1\,\% \\
\bottomrule
\end{tabularx}
\end{table}

\textit{Nota.} Los datos se obtuvieron mediante encuesta propia aplicada a 108 conductores de la ciudad de Ayacucho (Huamanga) durante el mes de julio de 2026.

\subsubsection{Problemática identificada en estacionamientos}

Se aplicó una encuesta complementaria a 21 trabajadores y administradores de estacionamientos:

\begin{table}[H]
\caption{Resultados de Encuesta a Trabajadores y Administradores}
\begin{tabularx}{\textwidth}{X r}
\toprule
\textbf{Indicador} & \textbf{Porcentaje} \\
\midrule
Considera un problema no saber cuántos espacios quedan disponibles & 66,7\,\% \\
Identifica congestión durante las horas punta & 57,1\,\% \\
No utiliza actualmente un sistema especializado de administración & 61,9\,\% \\
Está interesado en conocer la disponibilidad en tiempo real & 95,2\,\% \\
Considera importante automatizar el registro de ingreso y salida & 100\,\% \\
Utilizaría una plataforma para administrar el estacionamiento & 85,7\,\% \\
Implementaría un sistema tecnológico de gestión & 95,2\,\% \\
Pagaría una suscripción mensual por el servicio & 90,5\,\% \\
Considera importante ahorrar tiempo & 90,5\,\% \\
Considera importante tener mayor control del negocio & 85,7\,\% \\
\bottomrule
\end{tabularx}
\end{table}

\textit{Nota.} Los datos se obtuvieron mediante encuesta propia aplicada a 21 trabajadores y administradores de estacionamientos de la ciudad de Ayacucho (Huamanga) durante el mes de julio de 2026.

\subsubsection{Solución propuesta}

Ante esta problemática se propone Smart-Park, cuyas capacidades principales son:

\begin{itemize}
    \item Registro y administración centralizada de múltiples estacionamientos afiliados, mediante un flujo de solicitud de afiliación que el administrador de plataforma revisa, aprueba o rechaza.
    \item Localización de estacionamientos mediante geolocalización y mapa interactivo con capas de calles y satélite, y filtros por zona, tipo de espacio y tarifa.
    \item Consulta en tiempo real de la disponibilidad, las tarifas y las calificaciones, con actualización inmediata del estado de los espacios mediante WebSocket.
    \item Estudio de diseño de planos (editor CAD en dos dimensiones) para que cada local dibuje su distribución física: espacios para autos y motos, espacios techados, muros, accesos, garita y pasos peatonales.
    \item Reserva de un espacio específico sobre el plano interactivo o asignación automática del primer espacio disponible, en las modalidades estándar (por horas o minutos), anticipada (fecha futura) y abono (semanal, mensual o fraccionado por días).
    \item Pase de acceso digital con código QR y temporizador de tolerancia, verificable por el personal de garita en una página pública de validación.
    \item Liberación automática de reservas vencidas mediante un tiempo de tolerancia configurable por local.
    \item Cobro automático según el tiempo de permanencia, la tarifa del tipo de vehículo, la unidad de cobro y el recargo nocturno, con liquidación de la sobreestadía por la pasarela o en la garita.
    \item Pagos electrónicos con tarjetas, billeteras digitales (Yape y Plin) y PayPal, a través de las pasarelas Culqi y PayPal, con cálculo de comisiones y liquidación de fondos a cada local.
    \item Detección automática de vehículos mediante cámaras IP y visión artificial (modelo YOLOv8, con respaldo en OpenCV), habilitada en los locales que aprueban una evaluación técnica de techado, calidad de cámaras, estabilidad de red, trazado de espacios e iluminación.
    \item Seguridad basada en control de acceso por roles (RBAC), sesiones con tokens JWT en cookies HttpOnly, inicio de sesión con correo y contraseña o con cuenta de Google, PIN de seguridad de 4 a 6 dígitos para los perfiles administrativos, limitación de intentos de inicio de sesión y bitácora de auditoría.
    \item Notificaciones en tiempo real, reseñas y calificaciones de 1 a 5 estrellas, reporte de incidencias con evidencia fotográfica, y administración de términos, condiciones y políticas de uso.
    \item Reportes estadísticos y financieros a nivel de local y de plataforma, con copias de seguridad automáticas de la base de datos.
    \item Interfaz web responsiva, instalable como aplicación (PWA) en computadoras, tablets y teléfonos móviles.
\end{itemize}

\subsubsection{Actores del sistema}

\paragraph{Actores principales} Se identifican cuatro actores principales que interactúan directamente con el sistema:

\begin{table}[H]
\caption{Actores Principales del Sistema}
\begin{tabularx}{\textwidth}{p{2.8cm} p{2cm} p{2.5cm} X}
\toprule
\textbf{Actor} & \textbf{Tipo} & \textbf{Nivel} & \textbf{Responsabilidades principales} \\
\midrule
Administrador de plataforma & Principal (interno) & Nivel central & Administra toda la red de estacionamientos afiliados. Aprueba o rechaza afiliaciones, realiza la evaluación técnica y habilita el módulo de cámaras, gestiona usuarios y roles, define comisiones y configuración global, liquida fondos, modera reseñas, envía comunicados masivos y supervisa la auditoría. \\
\midrule
Administrador de local & Principal (interno) & Central y estacionamiento & Administra un estacionamiento específico: datos generales, plano digital, espacios, tarifas, políticas de reserva, cámaras y su calibración, personal y turnos. Responde reseñas, resuelve incidencias y consulta los reportes de su local. \\
\midrule
Trabajador del local & Principal (interno) & Estacionamiento & Supervisa la operación diaria: valida los pases QR, registra ingresos y salidas, controla el estado de los espacios, cobra en garita cuando corresponde y registra incidencias. \\
\midrule
Usuario final (conductor) & Principal (externo) & Nivel central & Se registra, busca estacionamientos en el mapa, registra sus vehículos, reserva y paga espacios, usa el pase QR, amplía o cancela reservas, consulta su historial, reporta incidencias y califica el servicio. \\
\bottomrule
\end{tabularx}
\end{table}

\paragraph{Actores secundarios y sistemas externos} Complementan las funcionalidades del sistema:

\begin{table}[H]
\caption{Actores Secundarios y Sistemas Externos}
\begin{tabularx}{\textwidth}{p{4.5cm} X}
\toprule
\textbf{Actor o sistema} & \textbf{Función en el sistema} \\
\midrule
Cámaras IP de los locales habilitados & Proporcionan las imágenes que el motor de visión artificial analiza para detectar vehículos y determinar la ocupación de los espacios. \\
Motor de visión artificial (YOLOv8 / OpenCV) & Detecta los vehículos (autos, motos, buses y camiones) y clasifica cada espacio como libre u ocupado. \\
Pasarelas de pago (Culqi y PayPal) & Procesan los cobros con tarjeta y billeteras digitales, y confirman o rechazan cada transacción. \\
Servicio de autenticación de Google & Permite iniciar sesión con una cuenta de Google. \\
Servicio de mapas (Leaflet) & Muestra el mapa interactivo, la ubicación de los locales y la posición del usuario. \\
Procesos automáticos del sistema & Liberan reservas por vencimiento de la tolerancia, escanean las cámaras de forma periódica y realizan las copias de seguridad programadas. \\
\bottomrule
\end{tabularx}
\end{table}

\subsection{Diagrama de Contexto}

El diagrama de contexto (Figura~1) sitúa al sistema Smart-Park como bloque central. Alrededor de este bloque se representan los actores y sistemas externos, cada uno conectado mediante una flecha etiquetada con el tipo de dato que intercambia.

% Aquí debes incluir tu imagen del diagrama de contexto:
% \begin{figure}[H]
% \caption{Diagrama de Contexto del Sistema Smart-Park}
% \centering
% \includegraphics[width=0.95\textwidth]{imagenes/diagrama_contexto.png}
% \end{figure}
% \textit{Nota.} Elaboración propia.

La Tabla~7 describe los flujos de información entre el sistema y cada actor externo:

\begin{table}[H]
\caption{Intercambio de Datos en el Diagrama de Contexto}
\begin{tabularx}{\textwidth}{p{3.5cm} p{2cm} X}
\toprule
\textbf{Actor / Sistema externo} & \textbf{Dirección} & \textbf{Datos intercambiados} \\
\midrule
Usuario final (conductor) & Bidireccional & Envía: ubicación, datos de reserva, placa y tipo de vehículo, pago, reseñas e incidencias. Recibe: disponibilidad en tiempo real, tarifas, confirmaciones, notificaciones, pase QR y comprobantes. \\
\midrule
Administrador de plataforma & Bidireccional & Envía: decisiones de afiliación, resultados de evaluación técnica, comisiones, configuración global y comunicados. Recibe: solicitudes de afiliación, reportes consolidados y alertas de auditoría. \\
\midrule
Administrador de local & Bidireccional & Envía: datos del local, plano, tarifas, horarios, políticas de reserva, personal y registro de cámaras. Recibe: reportes del local, reseñas e incidencias. \\
\midrule
Trabajador del local & Bidireccional & Envía: validación de pases QR, registro de ingreso y salida, e incidencias. Recibe: estado de los espacios en tiempo real y alertas operativas. \\
\midrule
Cámaras IP & Entrante & Imágenes (MJPEG o JPEG) de los locales habilitados, capturadas por la plataforma mediante HTTP/HTTPS. \\
\midrule
Pasarelas de pago & Bidireccional & Envía: solicitud de cobro. Recibe: confirmación o rechazo de la transacción. \\
\midrule
Autenticación de Google & Bidireccional & Envía: solicitud de inicio de sesión. Recibe: identidad verificada del usuario. \\
\midrule
Servicio de mapas & Bidireccional & Envía: coordenadas del usuario y de los locales. Recibe: capas cartográficas de calles y satélite. \\
\bottomrule
\end{tabularx}
\end{table}
